\section{Discussion}
\label{sec:discussion}

\subsection{Implications for reported results}
The seven corpora fall into three groups. PhiUSIIL and the PT-Alexa construction are separable by recording conventions, and the contrast persists among homepages alone (Table~\ref{tab:controls}): \texttt{https://www.} in front of every legitimate PhiUSIIL string, a missing scheme in front of every legitimate PT-Alexa string, and no path in any legitimate URL of either. On these corpora raw-string scores say little about phishing detection; canonical results are the more informative numbers, and they are lower by \DRawVsCanonPhiUSIIL{} and \DRawVsCanonPTAlexa{} MCC. Because their legitimate classes contain only homepages, even canonical results there partly measure the homepage contrast, which no transformation of the strings can remove.

Kaggle-MU is a different case. Its character model depends on formatting (a drop of \DRawVsCanonKaggleMU, \KCleanDrop{} on the cleaned subset and \KCleanTwoDrop{} without the bare-\texttt{www.} phishing rows), but the larger problem is its labels: thousands of canonical strings and registered domains with both labels, a standards body as its largest phishing domain, thousands of URLs that PhishStorm labels the other way, and thousands of benign-looking pages labelled phishing. Formatting and source mixing cannot be separated in this collection. It also contains \ISCXInKagglePct\% of ISCX-URL2016. It should not be used as an independent benchmark next to ISCX-URL2016 or PhishStorm. Its low transfer scores are consistent with these label inconsistencies, although we have not tested that they cause them.

Ebbu2017, PhishStorm, ISCX-URL2016 and Hannousse2021 have moderate presentation probes, but a character n-gram model loses less than 0.01 MCC on average when the formatting is removed (a CNN up to \CNNDropHannousseRandom{} under the random split), and grouping by registered domain costs the canonical n-gram model at most \DCanonRandomVsDomainPhishStorm{} MCC there, the lexical model up to \DLexicalRandomVsDomainISCXURL{} and the CNN up to \CNNGroupCostISCXURL{} on ISCX-URL2016. They are the most usable of the seven for URL-only detection, with one caveat: ISCX-URL2016 draws its legitimate class from \RegLegitISCXURL{} registered domains, so any detector that weights the host heavily needs a grouped split there.

For a security operations team the operating point matters more than the balanced score. Within a corpus, a threshold set on validation data for a false-positive rate of 0.1\% holds that rate on the test data and gives the canonical model a precision between \VPrecOneKaggleMUCharCanon{} and \VPrecOneISCXURLCharCanon{} at 1\% prevalence (Table~\ref{tab:precision}), at true-positive rates between \VTprKaggleMUCharCanon{} and \VTprISCXURLCharCanon. Across corpora the same procedure fails: the realised false-positive rate has a median of \XVFprCharCanon\% and exceeds ten times the target in \XVOverTenCharCanon{} of \XVNPairsCharCanon{} pairs, and the median precision at 1\% prevalence would be \XVPrecOneCharCanon, close to the base rate itself. A threshold therefore has to be set on traffic from the deployment, and a benchmark score does not tell an operator where to set it. A filter tuned on PhiUSIIL raw strings would also have learned that \texttt{https://www.} homepages are legitimate, which is the opposite of what a defender wants once phishing sites use HTTPS. These findings do not show that published detectors are weak. They show that reported numbers mix detection skill with corpus-specific conventions in unknown proportions, and that the proportions differ sharply between corpora.

What a deployed detector sees also differs from what these corpora store. A mail gateway or a browser extension sees full URLs with a scheme, often wrapped by a rewriting service or a click tracker that has to be unwrapped first; a DNS resolver, or a proxy without TLS interception, sees only the host name. For the second group the host-view results are the relevant ones, and they are the results most affected by domain leakage. Phishing hosted on legitimate platforms is a further limit of URL-only detection: \texttt{\LargestPhishPhiUSIIL}, \texttt{\LargestPhishPTAlexa} and \texttt{\LargestPhishISCXURL} are the registered domains with the most phishing URLs in PhiUSIIL, PT-Alexa and ISCX-URL2016, so a host-only model cannot separate those URLs from benign pages on the same platform.


\subsection{A reporting checklist for URL benchmarks}
Table~\ref{tab:checklist} turns the audit into checks that a benchmark paper can run with the audit code. None of them requires page content or access to the original collection pipeline, and the full set ran in \RunMinutes{} minutes for all seven corpora on four CPU cores. The last column gives the action we suggest when a check fails; the thresholds are judgement calls informed by the spread of our results, not statistical cut-offs.

\begin{table}[tp]
\centering
\caption{Checks recommended before reporting results on a URL corpus, with a suggested action when a check fails. The audit code implements C1 to C5, C7, C8 and the transfer part of C6; the time-ordered split of C6 is implemented only for the PhishTank dump, the one source with submission times.}
\label{tab:checklist}
\footnotesize
\begin{tabularx}{\linewidth}{@{}>{\raggedright\arraybackslash}p{2.7cm}>{\raggedright\arraybackslash}X>{\raggedright\arraybackslash}X@{}}
\toprule
Check & What to report & Action if it fails \\
\midrule
C1 Duplicates and containment & Exact and canonical duplicate counts, strings carrying both labels, and the share of the corpus contained in other benchmark corpora. & De-duplicate canonically before splitting; do not count contained corpora as independent evidence. \\
C2 Presentation probe & MCC of a classifier on presentation flags only, with and without the scheme type, and the per-flag ablation. & If the probe is high, run C3 and C4 before interpreting any raw-string score. \\
C3 Canonical input & Results on canonical strings next to raw-string results, with paired differences over repeated splits and a corrected test. & If the drop exceeds a few hundredths of MCC, report canonical results as the primary numbers. \\
C4 Sampling probe & Structure-only and host-only results, and the probe on path and homepage subsets. & If one class has no paths, say so and report homepage-subset results. \\
C5 Domain grouping & Share of test URLs whose registered domain occurs in training, and results under a stratified split grouped by registered domain. & Report grouped results as primary for host-weighted detectors. \\
C6 Transfer and time & ROC-AUC on at least one other corpus after purging test URLs against the full training corpus by canonical URL and by registered domain, and a time-ordered split where collection dates exist. & Check whether a good transfer result rests on a convention shared with the test corpus (C2 on both). \\
C7 Operating point & Precision at a realistic prevalence at the reported threshold and at a threshold set on validation data for a fixed low false-positive rate, with the realised test rate. & Report the threshold and how it was chosen. \\
C8 Provenance & Mirror or archive version, file digests and the upstream chain, because public copies of one corpus differ in presentation. & Name the exact copy used. \\
\bottomrule
\end{tabularx}
\end{table}

\subsection{Threats to validity}
\label{sec:threats}
\paragraph{Construct validity} The boundary between artefact and signal is not sharp. The \texttt{https} flag is partly a real property whose prevalence among phishing URLs rose over the years these corpora span (in the PhishTank dump behind PT-Alexa from \PTHttpsEarly\% of submissions up to 2017 to \PTHttpsTwenty\% in 2020), so a scheme-type contrast mixes collection year with collection pipeline; we report every probe with and without it. The PhishTank series describes the submitted URL, often the lure link of a message, not necessarily the landing page, and it is therefore the view of a mail gateway rather than of a browser. On PhiUSIIL the flag cannot be read as a server property at all, because every legitimate string carries the same \texttt{https://www.} prefix, which is a curator's template. A missing \texttt{www.} or scheme is not equivalent to its presence under the URI standard \citep{rfc3986}, and we claim only that in these corpora such differences mainly track the pipeline. Canonicalisation keeps userinfo, ports, path, query and fragment, so it removes no attacker-controlled content beyond the formatting flags, but a phishing campaign that prefers \texttt{http://} or omits \texttt{www.} would lose a weak genuine cue. Structure features are not pure artefact either, since phishing pages do sit deeper in sites, and the fragment we keep is partly a pipeline property too, because crawls rarely record fragments while feeds store what reporters pasted. The canonicaliser drops ports 80 and 443 whatever the scheme, treats a defanged scheme such as \texttt{hxxp://} as a scheme, and does not normalise hexadecimal or decimal IP hosts. Internationalised host names are not converted between Unicode and punycode, so a feed that stores one form and a crawl that stores the other would differ in a property we do not audit.

\paragraph{Internal validity} No corpus in this study was collected by one pipeline for both classes, so we have no positive control in which the probe should fall to chance. The path and homepage subsets are the closest substitutes: they remove the sampling contrast but not the collection pipeline. Corpora are capped at 25,000 URLs per class, which changes the class ratio of several corpora (Table~\ref{tab:corpora} gives the original ratios); for Hannousse2021 and PT-Alexa the cap exceeds the class sizes, so their ten seeds vary only the split. The registered-domain grouping relies on a Public Suffix List snapshot, and URL shorteners or hosting providers outside its private section can place unrelated URLs in one group, which makes the grouped split conservative. The CNN is a single architecture trained with three seeds, fixed hyper-parameters and a fixed number of epochs. The Kaggle-MU-clean subset removes rows by rules we chose after seeing the conflicts, so it is a sensitivity check, not a corrected corpus.

\paragraph{External validity} We audit public mirrors, not the original releases, and copies of one corpus differ in the properties we measure (Section~\ref{sec:corpora}); results for a corpus hold for the named mirror and commit. Phishing feeds themselves contain entries that are not phishing pages \citep{marchal2018}, which adds label noise we cannot measure. All corpora are balanced or close to balanced, while deployment base rates are far lower \citep{marchal2018,dalton2025}; the precision conversion in Section~\ref{sec:canon} is an estimate that assumes the false-positive rate of the test split carries over to live traffic. The corpora span 2014 to 2024. We model time through the publication-year ordering of transfer pairs, a proxy that misplaces corpora such as Kaggle-MU that reuse older collections, and through one time split of the PhishTank dump, the only source with submission times; collection windows inside the other corpora are mostly undocumented. Phishing URLs are short-lived and blocklists react within hours to days \citep{oest2020phishtime,oest2020sunrise}, so a snapshot of a feed can over-represent campaigns that stayed online long enough to be reported and verified.
